\section{Full-cover movement and its ordinary boundary}
\label{app:movement}
\label{app:swap}
The finite-cover and affine-routing method develops the constructions
of~\cite{swap,mixed-precompile}; the full-weight normalization and
ordinary boundary below are proved for the present rational body.

This appendix proves the movement interface of Lemma~\ref{lem:movement}
from the finite template of Lemma~\ref{lem:body}. An \emph{atom} is an
address field of $w$ base-$q$ digits, of range $q^w$. A coarse field
consists of several such atoms; the payload and every other coordinate
remain spectators during an ordinary interchange. If a call has active input array length $V$, all costs below
are divided by this actual $V$. In particular $V$ includes retained pads
and opaque payloads in the current call. Cells in a parked ancestor frame
are not part of the active call and are not traversed by its descendants;
the parent's scans that park, restore, and erase such a frame are charged
to the parent's own $O_N(V)$ node toll.

\subsection{The rational address body and its full-cover endpoint}

Fix the finite body supplied by Lemma~\ref{lem:body}, with rank $N\ge2$,
$B$ independent productive arrays, and $P$ paid rank-one queries. Put
\[
 a=P/B.
\]
This is a real number: cover rows and role multiplicities may give
fractional normalized coefficients. A rank-one query sees the whole
current record array divided by $B$, hence has volume exactly $V/B$;
this includes owned queries. Let $z_0,z$ be fixed rationals with
\[
 s_0=1-\frac1{1300}<z_0<z<1.
\]
The body lemma allows $N$ to be chosen so that $a<\varepsilon N^z$
for any prescribed fixed $\varepsilon>0$.

The address action is over a rational vector space; the payload action is
over $\mathbb F_2$. These are separate rings. Write $R$ for the fixed
two-dimensional rational matrix of order three used in the address
representation. The frozen native body realizes a full address operator
$M_g$ on every cover row $g$, with
\begin{equation}
 M_g^2+M_g+I=0,\qquad M_g^3=I.                 \label{eq:movement-1}
\end{equation}
The equality is an identity of address matrices and is used only for
composition of address maps. The corresponding permutation of payload
arrays is not asserted to satisfy an additive polynomial identity.
All adapters and native jobs are executed on complete role arrays, with
arbitrary independent initial contents.

After freezing the finite body and its finite rational matcher list,
choose a prime $q\ge5$ avoiding every denominator and every nonzero
witness coefficient occurring below. For $w\ge1$, let
\[
 Q=q^w,\qquad \mathcal G_w=\mathrm{GL}_{2N}(\mathbb Z/q^w\mathbb Z),
 \qquad d=2N.
\]
Each cover row is indexed by $g\in\mathcal G_w$ and carries its own
$M_g=gTg^{-1}$, where $T$ is the frozen full order-three address matrix.
The cover is the full group, not a subgroup presumed to contain the
rational matchers. Every fixed matcher acts by a permutation of this
cover.

For a paid leaf, the matrix $M-I$ has split rank two over $\mathbb Q$.
On a good chart, unit-lower row and column operations normalize its
selected cross block to $ED F^{\mathsf T}$, where $E,F$ select two
coordinates and $D$ is diagonal with unit entries. Write
\[
 M-I=UV,\quad U=\binom{E}{B_0},\quad V=(A_0,DF^{\mathsf T}),
 \quad H=VU=A_0E+DF^{\mathsf T}B_0,
\]
and set
\[
 K=(-B_0+FD^{-1}H)E^{\mathsf T}
       +FD^{-1}A_0(I-EE^{\mathsf T}).
\]
The selector identities $E^{\mathsf T}E=F^{\mathsf T}F=I_2$ give,
by direct multiplication,
\[
 KE=-B_0+FD^{-1}H,\qquad DF^{\mathsf T}K=A_0,
\]
and therefore, with $L(K)=\left(\begin{smallmatrix}I&0\\K&I\end{smallmatrix}\right)$,
\begin{equation}
 L(K)ML(-K)=I+
 \binom{E}{FD^{-1}H}(0,DF^{\mathsf T}).             \label{eq:movement-2}
\end{equation}
The residual in \eqref{eq:movement-2} is supported on four coarse fields, regardless
of $N$. Only the selected unit pivots are inverted; $H$ is not inverted.
Thus four-field Gaussian lowering applies over the local ring
$\mathbb Z/q^w\mathbb Z$.

We record the chart count. For each leaf, require its selected northeast
one-by-one and two-by-two minors to be units. A witness for their
simultaneous nonvanishing is obtained by taking $E$ to be the first two
head columns, $F$ the last two tail columns in reverse order, and
$H_0=R-I_2$; with $U=(E^{\mathsf T},(FH_0)^{\mathsf T})^{\mathsf T}$ and
$V=(0,F^{\mathsf T})$, one has $VU=H_0$, $\det H_0=3$, and
$I+UV$ is rationally similar to $R\oplus I$. The selected minors are
nonzero (indeed units after the fixed conjugacy denominators are avoided).
For the full output require the entire $N$-by-$N$ upper-right block to
be invertible; the witness $R\otimes I_N$ has upper-right block
$-I_N$.

After clearing the fixed conjugacy denominators, each entry of
$gT\operatorname{adj}(g)$ has degree at most $d$. The product of the
two selected leaf-minor polynomials has degree at most $3d=6N$; the
full-output determinant has degree at most $Nd=2N^2$. The number of
invertible matrices in $\mathrm{GL}_d(\mathbb F_q)$ is at least
$3/4$ of all matrices, since
\[
 \prod_{i\ge1}(1-q^{-i})\ge1-\sum_{i\ge1}q^{-i}\ge3/4.
\]
The polynomial zero bound, followed by conditioning on invertibility,
therefore gives bad fractions at most $12N/q$ for a leaf and $4N^2/q$
for the full output. We use the common bound
\begin{equation}
 \delta_{\rm bad}=16N^2/q.                           \label{eq:movement-3}
\end{equation}
Every low matrix $g_0\in\mathrm{GL}_d(\mathbb F_q)$ has exactly
$q^{d^2(w-1)}$ lifts, so these same fractions apply at every $w$.
The average is taken separately for each leaf type and for the root
output type; no union bound over occurrences in the finite word is used.

On a good leaf, unit-pivot Gaussian elimination on four fields costs at
most
\[
 4+2\cdot4\cdot3+5\cdot3=43< G:=64
\]
ordinary interchanges. Here an upper addition is one lower addition
conjugated by two swaps. A paired unit scaling uses five swaps by
\[
 J_a=\operatorname{NegLast}\,U(a)L(-a^{-1})U(a),\qquad
U(a)=J_1L(a)J_1,
\]
followed by $J_1$; fixed signs and lower additions are streamed.
On a bad leaf or bad root row, Gaussian elimination on all $2N$ fields
uses at most $F_N:=32N^2$ ordinary interchanges. No nonunit is inverted.
These interchanges are the only recursive calls; all other row/column
operations in the fixed normalizers are streamed affine updates.

\subsection{Removing the full matrix weight}

The good-root adapters must remove the variable matrix weight of the full
cover, not merely a canonical representative. Write
\[
 M=\begin{pmatrix}X&B_0\\C&D_0\end{pmatrix},
 \qquad M^2+M+I=0,
\]
and assume $B_0$ is invertible. The block identities imply
$D_0=-I-B_0^{-1}XB_0$ and direct multiplication yields
\[
 L(B_0^{-1}X)M L(-B_0^{-1}X)
   =\begin{pmatrix}0&B_0\\-B_0^{-1}&-I\end{pmatrix}.
\]
After a lower adapter and a head negation this is
\[
 J_A=\begin{pmatrix}0&A\\A^{-1}&0\end{pmatrix},\qquad A=-B_0.
\]
For this full, possibly variable $A$, block multiplication gives
\begin{equation}
 J_A L(-A^{-2})J_A=U(-I),\qquad
 J_I=\operatorname{NegHead}L(I)U(-I)L(I).             \label{eq:movement-4}
\end{equation}
Consequently two completed calls to the same body, with the lower
adapters in \eqref{eq:movement-4}, cancel the full weight and leave the prescribed
canonical endpoint. The body restores all cover and role fields before
returning, so the second call sees the same $g$ and the same $A(g)$.
No variable unpaired scaling is delegated to a child.

On bad root rows, set every wrapper adapter to the identity, but still
execute both complete bodies on every row. Their endpoint there is
$M_g^2$. After the completed calls, apply the fixed correction $J_IM_g$
by the $F_N$-interchange fallback; since $M_g^3=I$, this gives $J_I$.
This branch is taken only at completed endpoints, because intermediate
right routes can mix row labels.

The normalized child coefficient is therefore bounded by
\begin{equation}
 \mu_z\le {2Ga+\delta_{\rm bad}F_N(2a+1)\over N^z}.       \label{eq:movement-5}
\end{equation}
The $2Ga$ term counts the two bodies on every row; the bad-leaf fallback
budgets are included in the two $a$ contributions, and the final
$\delta_{\rm bad}F_N$ is the one completed-root correction. The extra
factor in $(2a+1)$ is a conservative allowance for the full bad-row
branches while preserving the same volume normalization. More precisely,
coefficients in \eqref{eq:movement-5} are real normalized totals: the fraction of rows
in a chart class multiplies its number of child calls, while each query
volume is $V/B$. This is not an integer number of calls per individual
row. Owned queries are included in $P$ and therefore in $a$.

Choose $N$ using Lemma~\ref{lem:body} with $\varepsilon=1/512$, so
$a<N^z/512$. Then $2Ga/N^z<1/4$. Choose $q$ after the finite body and
all witnesses are frozen, large enough for \eqref{eq:movement-3}, prime avoidance, and
\[
 q>2048N^4(2P+B).
\]
Using $N^z\ge1$, $B\ge1$, and $a=P/B$, this implies
$\delta_{\rm bad}F_N(2a+1)/N^z<1/4$. Thus $\mu_z<1/2$; in particular
$\mu_z<3/4$. Since $N>1$ and $z<1$, $\mu_1\le\mu_z<1$.

\subsection{Fixed rational routes and batched descriptors}

For every fixed rational matcher $S$ in the frozen word, all denominators
are units modulo $q^w$. Write a cover index as
\[
 g=g_0+qG_1,\qquad g_0\in\mathrm{GL}_{d}(\mathbb F_q),
 \quad G_1\in(\mathbb Z/q^{w-1}\mathbb Z)^{d^2}.
\]
The low type $g_0$ ranges over a fixed finite set. On each such type,
right multiplication by $S$ induces a fixed affine bijection
\begin{equation}
 G_1\longmapsto G_1S+c(g_0)\pmod {q^{w-1}}.             \label{eq:movement-6}
\end{equation}
The fixed linear part in \eqref{eq:movement-6} factors into a fixed number (depending
only on $N,q,S$) of coordinate swaps, shears and fixed-unit scalings.
Lower shears are rotations of complete target fibers. Upper shears use
the previously compiled ordinary interchange on fields of width $O(w)$.
Fixed scalings use a fixed number of residue lanes. These route operations
are independent of input data and of the cover index except through its
streamed descriptor; variable Gaussian pivots are still charged by the
child calls counted in \eqref{eq:movement-5}. No finite subgroup of $\mathrm{GL}_d(\mathbb Q)$
is assumed.

Here is the array shape that supplies the target fibers. At recursion
depth $j$, split the unused numeric row stock as
\[
 [R_j]\times[B]\times[|\mathrm{GL}_d(\F_q)|]
 \times\prod_{i=1}^{d^2}[q^{w-1}].
\]
After this row factor the array contains two complete $n$-atom address
fields, an arbitrary intervening range, spectator coordinates and an
opaque bit payload. These are
Cartesian factors of its lexicographic index, not headers appended to
each payload record. The $B$ factor selects independent source streams;
the next factors index $g_0$ and the $d^2$ entries of $G_1$. Thus fixing
a chart, a role, every spectator and all but any one target atom leaves
the complete target fiber $[q^w]$. Good and bad chart conditions depend
only on $g_0$, not on that target coordinate. Owned arrays have the same
index set even when their initialized values are zero.

Right routing permutes only the chart factors and preserves this product
shape at completion. A normalizer updates a target atom by a bijection
depending on the other factors, so its completed array again has every
target position. At a recursive call the consumed $B$ and chart factors
are fixed while the independent next row factor and all smaller target
atoms remain complete. Induction proves the target-fiber property at
every depth; it is independent of the payload values.

Restriction to a good or bad chart class fixes only the current low type
$g_0$. Packing those rows in their inherited order preserves $R_j$ as an
independent complete Cartesian factor. Right routing acts on the current
$g$ factor alone. Hence the next $W(w)$ factor extracted from $R_j$ remains
uniform conditional on every ancestor chart choice. Current cover labels
remain spectators and are not divided from query volume a second time.
The scans extracting and restoring a chart class cost $O_N(V)$ at its
parent and are included in that parent's node toll.

An odometer enumerates the chart indices. Hold one current fiber header
on work tapes and generate its offsets before scanning that fiber; do
not copy chart digit strings into each record. The finite $g_0$ table,
fixed-size rational matrices and low-type pivot choices are fixed data.
Matrix arithmetic modulo $q^w$ uses schoolbook operations and unit
inversion with polynomial bit cost in $w$. Ancestor headers are stack
entries on a fixed tape set. Their depth is $O_N(\log(s+2))$; on the
chosen root families $s=O(w^{1/\delta_0})$, so even redundant sequential
scans and recomputation cost $C_N(w+1)^r$ for a fixed integer $r$.
This also bounds normalizer coefficients and offset endpoints. They
are generated once per complete target fiber, not once per payload bit.
There is one batched recursive invocation, not a separately allocated
program per high matrix. Since each fiber has $q^w$ positions, for every
fixed integer $r$, there is an effective $C_{q,r}$ with
\begin{equation}
 (w+1)^r\le C_{q,r}q^w\qquad(w\ge1),                \label{eq:movement-7}
\end{equation}
so fixed-degree descriptor work is absorbed by that fiber's record scan.
All counters use fixed-width digit lanes with ripple carries; resets and
head returns are charged to the full fiber scan. The stack depth is at
most $D=O_N(\log(s+2))$, and its descriptor volume is
$O_N(w\log(s+2))$ digits. This is the only growing stack dimension and
it is represented on a fixed number of tapes with separators; no tape is
allocated per chart, axis, or recursion depth.

For completeness, multiplication by a fixed integer unit also has an
explicit streaming implementation. Let $Q=q^w$ and let $a>0$ be fixed
with $(a,q)=1$. Partition $x\in[0,Q)$ into the at most $a$ intervals on
which $j=\lfloor ax/Q\rfloor$ is constant. On the $j$th interval the
image $y=ax-jQ$ is increasing and satisfies $y\equiv-jQ\pmod a$.
Since $(Q,a)=1$, these residues are distinct. A fixed $a$-lane merge
therefore emits the records in increasing $y$ order, implementing
$x\mapsto ax\pmod Q$ in $O_a(Q)$ record moves. For the inverse, partition
the input $y$ by the residue
$r_j=-jQ\pmod a$; in that lane set $x=(y+jQ)/a$. Its image interval is
$[jQ/a,(j+1)Q/a)$, so concatenating the $a$ lanes in order emits every
$x\in[0,Q)$ once, again in $O_a(Q)$ moves. Numerator, denominator and
sign lanes compose to implement any fixed rational unit route. The lane
count depends only on the frozen rational coefficient. This construction
does not implement a variable unit for free.

\subsection{The stopped recurrence and actual row stock}

Let $b\le1$ be the exponent of a previously compiled ordinary
interchange supplier. Fix an admissible root width $s$ and atom width
$w=\lceil s^{\delta_0}\rceil$, with $0<\delta_0<1$. Let $T(n,w)$ be
the total bit time divided by the actual active volume for the covered
body on $n$ atoms per field, averaged with the real row weights described
after \eqref{eq:movement-5}. At a node, put $f=\lfloor n/N\rfloor$. The recursive children
are calls to this same covered-body construction on strictly smaller
$n$; the predecessor supplier is used only for fixed routes, base cases,
and the fewer than $N$ remainder atoms. There are finitely many fixed
tape families: input/output, current roles, owned work, repair/padding,
and descriptor/parking stacks. Before a child, the parent parks only its
current sibling records and copies the child's active stream into an
adjacent frame. The child cannot scan parked ancestors. On return, the
parent restores that sibling and erases the frame. With fixed $N$ and a
fixed number of child types, all such parent operations cost $C_Nn$ per
node, not an additional multiplier on the descendant time.

The number of fresh cover/role labels per depth is
\[
 W(w)=B\,|\mathrm{GL}_{2N}(\mathbb F_q)|q^{4N^2(w-1)}.
\]
Take $D=\lceil\log_N(s+1)\rceil+1$ and choose $k$ with
$q^k\ge W(w)^D$. Reserve $k$ high q-adic digits in the head row field.
Pad that row field once to a multiple of $W(w)^D$; the expansion factor
is less than two. Every descent consumes one fresh $W(w)$ factor and
restores its labels only at the completed endpoint. The factor is divided
from the parent stock, not recreated in a child. Since $N,q,B$ are fixed,
\[
 k=O_N(w\log(s+2))=o(s),\qquad
 n=\left\lfloor {s-k\over w}\right\rfloor.
\]
The tail high field and all non-head fields remain spectators during the
body. After the complete body fixes the reserved head row, crop artificial
row padding. Use the predecessor supplier on the high head/tail pieces
(width $k$) and the low remainder pieces (width less than $w$); only after
these whole-field interchanges crop the original binary box. The initial
ordinal padding of each full binary address range to the least $q^s\ge2^e$
has factor less than $q^2$, a fixed constant.

The completed-call recurrence is
\begin{equation}
 T(n,w)\le \sum_{j}c_j T(\lfloor n/N\rfloor,w)
             +C_Nn+C_N(w+1)^b\quad(n\ge N),       \label{eq:movement-8}
\end{equation}
\begin{equation}
T(n,w)\le C_N(w+1)^b\quad(1\le n<N).             \label{eq:movement-9}
\end{equation}
The $c_j\ge0$ are real normalized multiplicities, not integer call
counts; their sum at exponent $z$ is bounded by \eqref{eq:movement-5}. Writing
$\mu_1=\sum_jc_j/N$ and $\mu_z=\sum_jc_j/N^z$, we have
$\mu_1,\mu_z<1$. Strong induction on $n$ proves
\begin{equation}
 T(n,w)=O_N\bigl(n+n^z(w+1)^b\bigr).               \label{eq:movement-10}
\end{equation}
Indeed, if $T(m,w)\le A m+B_0m^z(w+1)^b$ for smaller $m$, the child
sum in \eqref{eq:movement-8} is at most
$A\mu_1 n+B_0\mu_z n^z(w+1)^b$. Choose
$A\ge C_N/(1-\mu_1)$ and $B_0\ge C_N/(1-\mu_z)$, enlarging both to
cover \eqref{eq:movement-9}. The same induction applies to the real weighted sum because
all coefficients are nonnegative.

Substituting $n=O(e^{1-\delta_0})$ and $w=O(e^{\delta_0})$ into
\eqref{eq:movement-10}, and charging the high/low cleanup with the predecessor, gives
\begin{equation}
 O\!\left(V\left[e^{1-\delta_0}
 +e^{(1-\delta_0)z+\delta_0b}
 +e^{\delta_0}\log(e+2)+1\right]\right).          \label{eq:movement-11}
\end{equation}
The exponents in \eqref{eq:movement-11} are kept separately: no absorption is asserted
without comparing them. Set $\delta_0=1-z$. For $b>z$, the middle
exponent is
\[
 (1-\delta_0)z+\delta_0b
   =z^2+(1-z)b=z+(1-z)(b-z)>z,
\]
whereas the first exponent is $z$ and the reserve exponent is $1-z<z$.
The logarithm in the reserve term is absorbed in the positive gap
between $1-z$ and $z$. Thus the exact successor exponent is
\[
 b'=z+(1-z)(b-z),\qquad b_0=1,
\]
It follows that $b_j=z+(1-z)^{j+1}$. For any fixed $t>z$, a finite $j$ has
$b_j<t$. Compile stages in order. Stage $j$ uses stage $j-1$ only for
fixed routes, base/remainder work, and ordinary interchanges inside a
covered node. The covered children themselves use the same stage-$j$
construction at strictly smaller $n$, as in \eqref{eq:movement-8}; well-founded
induction on $n$ defines that recursion. There is no call to stage $j$ at
the atom-width parameter $w$. Therefore
the number of stages, tapes and all constants are fixed before the input
width varies, and this is a noncircular finite bootstrap. A zero-width
interchange is the identity and has a direct constant-per-input-bit
implementation, outside the positive-$n$ recurrence.
